% !PRINTED 165
\weberpage{165}
%<<CONT
durch $w$ theilbar ist, und dass, wenn $U$, $U'$, $U''\ldots$ durch $w$
theilbare, $V$, $V'$, $V''\ldots$ beliebige ganze rationale Functionen sind,
auch $UV + U'V' + U''V''\ldots$ durch $w$ theilbar ist.

Ist $W$ durch $w$ theilbar, so sagt man auch, $w$ geht in $W$ auf.

Zwei ganze rationale Functionen $U$, $V$, die nicht durch eine und dieselbe
ganze rationale Function theilbar sind, heissen \gesperrt{relativ prim}
oder \gesperrt{theilerfremd}.

Wir beweisen den Satz:

I. \gesperrt{Sind $U$, $V$, $v$ ganze rationale Functionen irgend welcher
Veränderlichen, sind $U$ und $v$ relativ prim und $UV$ durch $v$ theilbar,
so ist $V$ durch $v$ theilbar.}

Dieser Satz entspricht genau einem bekannten Fundamentalsatz aus der Lehre
von den ganzen Zahlen, dass nämlich, wenn ein Product von zwei ganzen
Zahlen durch eine dritte ganze Zahl theilbar ist, die zu dem einen Factor
theilerfremd ist, der andere Factor durch diese Zahl theilbar sein muss.

Wir beweisen ihn durch vollständige Induction. Sind $U$, $V$, $v$ nur von
\gesperrt{einer} Veränderlichen $x$ abhängig, so ist der Satz richtig; denn
nach \S.\,6 (Satz I) kann man in diesem Falle, wenn $U$ und $v$ relativ
prim sind, zwei andere ganze Functionen $P$ und $p$ von $x$ so bestimmen,
dass
\[
PU + pv = 1,
\]
woraus durch Multiplication mit $V$
\[
PUV + pVv = V
\]
folgt, und daraus ersieht man, dass, wenn $UV$ durch $v$ theilbar ist,
auch $V$ durch $v$ theilbar sein muss.

Wir nehmen also an, der Satz, den wir beweisen wollen, sei für Functionen
von $n$ und weniger Veränderlichen bewiesen, und wir leiten daraus seine
Richtigkeit für Functionen von $n + 1$ Veränderlichen ab.

Dazu ist erforderlich, dass wir aus dem als richtig vorausgesetzten
Theorem I einige Folgerungen ziehen, als deren Schluss sich dann die
Gültigkeit des Theorems für die nächst höhere Variablenzahl ergiebt.

Ist $v$ eine unzerlegbare Function, so ist eine andere Function $U$
derselben Veränderlichen entweder relativ prim zu $v$ oder durch $v$
theilbar. Daraus folgt nach I., dass ein Product von
%CONT>>
